\begin{table}[H]
\centering\centering
\caption{Race gaps in contract terms and across market margins \label{tab:pay-gap-terms}}
\centering
\resizebox{\ifdim\width>\linewidth\linewidth\else\width\fi}{!}{
\begin{tabular}[t]{lccccccc}
\toprule
  & (1) Guarantee share & (2) Zero guarantee & (3) Log years & (4) Log value & (5) Log APY & (6) Log APY & (7) Near minimum\\
\midrule
P(Black) & -0.012 & 0.015 & -0.041 & -0.038 & 0.024 &  & -0.004\\
 & (0.019) & (0.020) & (0.035) & (0.081) & (0.042) &  & (0.010)\\
P(other race) & -0.013 & 0.043 & -0.076 & -0.099 & 0.004 & 0.006 & -0.016\\
 & (0.034) & (0.037) & (0.063) & (0.135) & (0.070) & (0.071) & (0.015)\\
P(Black) $\times$ UFA, new team &  &  &  &  &  & 0.006 & \\
 &  &  &  &  &  & (0.053) & \\
P(Black) $\times$ UFA, re-signs with own team &  &  &  &  &  & 0.003 & \\
 &  &  &  &  &  & (0.054) & \\
P(Black) $\times$ extension &  &  &  &  &  & 0.079 & \\
 &  &  &  &  &  & (0.066) & \\
P(Black) $\times$ tag or tender &  &  &  &  &  & 0.044 & \\
 &  &  &  &  &  & (0.049) & \\
\midrule
Observations & 4686 & 5062 & 5062 & 5062 & 6511 & 6511 & 29634\\
R$^2$ & 0.353 & 0.234 & 0.516 & 0.656 & 0.756 & 0.756 & 0.530\\
Within R$^2$ & 0.240 & 0.192 & 0.451 & 0.629 & 0.732 & 0.732 & 0.518\\
Sample & Bargained, known guar. & Bargained & Bargained & Bargained & All veteran & All veteran & Non-rookie 2017+\\
Mean of outcome & 0.433 & 0.074 & 0.524 & 1.679 & 0.957 & 0.957 & 0.760\\
Implied Black gap (\%) &  &  & -4.0 & -3.7 & 2.4 &  & \\
Residual df & 4,009 & 4,386 & 4,386 & 4,386 & 5,829 & 5,824 & 28,991\\
\bottomrule
\multicolumn{8}{l}{\rule{0pt}{1em}* p $<$ 0.1, ** p $<$ 0.05, *** p $<$ 0.01}\\
\end{tabular}}
\end{table}

{\footnotesize
\noindent\textit{Notes:} This table includes the estimation results of equation (1) with the column-(5) controls of Table \ref{tab:pay-gap-veteran} (blocks A-D) and OTC position $\times$ year-signed fixed effects. Columns (1)-(4) use the main sample of freely bargained veteran contracts (unrestricted free agents and extensions) signed 2014-2026. Column (1): guaranteed money over total value, excluding contracts with zero recorded guarantees; a few contracts record guarantees above total value. Column (2): linear probability model of zero recorded guarantees (OTC records unknown guarantees as zero, so this mixes zero and unknown guarantees); it shows whether the sample of column (1) is selected by race. Column (3): log contract length in years. Column (4): log total contract value. Columns (5)-(6) use the full veteran market, which adds franchise and transition tags and RFA and ERFA tenders (pay set by a CBA formula; ERFA tenders are minimum deals); the outcome is log APY and the market-margin dummies are re-sign/extension and tag/tender (reference: UFA). Column (6) replaces P(Black) with its interactions with margin indicators, with a dummy for UFAs who sign with a team other than their end-of-season team of the previous year, so each coefficient is the gap within that margin with common control slopes. Expected Black contracts (sums of P(Black)) by margin: UFA new team 1539, UFA own team 1074, extension 476, tag or tender 831. Column (7) is a linear probability model of signing within 10\% of the reference minimum APY among non-rookie contracts (veteran market plus other, street free agent and practice-squad deals) signed 2017-2026; its market-margin dummies include the other/SFA/practice margin. Every column also includes fixed effects for the covariates of the race prediction's prior (position at NFL entry, entry era, draft-round bucket, college type and whether a home county is known), as regression calibration requires. Career stage (A): experience-bin dummies (0, 1, 2, 3, 4-6, 7+ seasons), age at signing and its square, a left-censored-experience indicator and an extension dummy (reference: unrestricted free agents). Whether a team extends a player is itself a team decision. Prior-season production (B): season-before-signing box-score and PFR measures with position-group-specific slopes (passing for QBs, rushing and receiving for RBs, receiving for WRs and TEs, tackles, sacks, QB hits, tackles for loss and pressures for DL and LB, tackles, interceptions, passes defended and coverage allowed for DBs, kicking and punting), games played and injury weeks interacted with position group, and a no-prior-season indicator. Career production (C): career games and injury weeks interacted with position group, and career sums of the block B production measures with the same position-group slopes. Pre-NFL signals (D): log draft pick, an undrafted indicator, draft-round dummies, 247 rating and stars, combine measures (40, vertical, bench, broad jump, cone, shuttle, height, weight) and athletic scores (a composite and size, speed, explosion and agility subscores built from the combine measures) interacted with position group, and the final college team's Power-conference, SRS and HBCU indicators. Controls with incomplete coverage are set to zero with a missing indicator (interacted like the control). The identifying assumption is that, conditional on the controls and fixed effects, race is uncorrelated with unobserved productivity. Omitted quality would bias the estimate: there are no PFF-type grades, and offensive linemen have no production measure, so their quality rests on games, injury weeks and pre-NFL signals. Regression calibration identifies the Black-white gap only if, in addition to names and hometown being unrelated to pay given race and the controls, (i) P(Black) is calibrated given every control in the column, not only given the prior's covariates, and (ii) the gap does not vary with the controls (otherwise the coefficient is a variance-weighted average of gaps). Controls that predict P(Black) within the prior's cells do not by themselves violate (i): under (i) they predict P(Black) exactly when they predict race. They do reduce the variation in P(Black) that identifies the coefficient (overlap), and if (i) fails for them the coefficient moves as they are added for reasons unrelated to pay. Table \ref{tab:pay-gap-veteran} reports the identifying variation by column and checks (i) against documented race. P(Black) is the predicted probability of being non-Hispanic Black alone; Black Hispanic and multiracial Black players enter P(other race), and Table \ref{tab:pay-gap-race-measures} also reports estimates with P(Black) + P(multiracial) as the Black regressor. The predicted Black share is not calibrated in levels: it lies below the TIDES African-American player share through 2016 and above the self-identified Black-alone share (below the Black-or-multiracial share) from 2019 (Table \ref{tab:race-prediction-tides}); Table \ref{tab:pay-gap-race-measures} reports estimates with the TIDES-raked posterior. Standard errors treat the predicted probabilities as known; the prior is estimated by EM on the whole player population of the database, so the omitted first-stage variance is likely small. The implied gap is $100(e^{\hat\beta_1}-1)$ for log outcomes. Residual degrees of freedom are observations minus estimated coefficients and fixed effects. Standard errors are clustered at the player level. P(Black) and P(other race) enter together, so the coefficient on P(Black) compares a Black (predicted: non-Hispanic Black alone) with a white player. Race is predicted, not observed: each person's probability of being non-Hispanic Black alone combines first name, surname and hometown county (BIFSG) with an NFL-specific prior estimated by EM on predetermined characteristics (players: position at entry, rookie era, draft round, college type, county availability; staff: role, unit and era at first appearance); documented race statements are not used (notes/race-prediction-design.md). Black Hispanic and multiracial Black persons are in the other-race probability, so P(Black) is narrower than the Black-alone-or-in-combination definition of the hand codes and of TIDES. Regressions use the probabilities (regression calibration) and control for the prior's covariates; the coefficient on P(Black) is the Black-white gap if the probabilities are calibrated given the regression's controls and names and hometown are unrelated to the outcome given race and the controls. The validation exhibit compares the predictions with published TIDES shares.
\par}

